\section{L'apprentissage pratique en ingénierie}\label{lapprentissage-pratique-en-ingenierie}

\stanceinfo{\textbf{FCEG AG de l'automne 2024} [2024-10-27]}{2024-10-27}

\officialstance{La Fédération canadienne étudiante de génie estime que l'enseignement de l'ingénierie doit comprendre un apprentissage pratique et engageant afin de préparer les futurs ingénieurs à leur future carrière. La Fédération canadienne étudiante de génie estime en outre qu'il est essentiel d'exposer les étudiants à l'apprentissage pratique dans le cadre des cours de base tout au long du diplôme et que les établissements devraient tenir compte de l'avis des étudiants sur l'apprentissage pratique.}

\stanceheading{La position des étudiants}
\begin{itemize}
 \item Le niveau d'apprentissage pratique actuellement disponible pour les étudiants canadiens en ingénierie dépend fortement de l'école qu'ils fréquentent et de la province dans laquelle elle est située.
 \item Bien que les étudiants puissent acquérir un apprentissage pratique par le biais de stages et de programmes coopératifs, ces opportunités ne sont pas garanties et ne sont pas exhaustives. Il serait donc bénéfique d'augmenter la disponibilité de l'apprentissage pratique offert tout au long des semestres universitaires en rendant les cours plus pratiques et en introduisant des environnements désignés dans lesquels les étudiants peuvent développer des compétences pratiques.
 \item En 2023, le BCAPG a mis à jour ses \href{https://engineerscanada.ca/sites/default/files/2023-12/Accreditation_Criteria_Procedures_2023.pdf}{critères et procédures d'accréditation} pour ajouter un nouveau critère de 3.4.4.4 comme indiqué ci-dessous :
 \begin{itemize}
  \item Au moins 225 unités d'accréditation du contenu du programme d'études en conception technique d'un programme d'ingénierie doivent être dispensées par des membres du corps professoral détenant un permis d'exercice du génie, comme le précise l'énoncé interprétatif sur les attentes et les exigences en matière de permis d'exercice [CITE].
  \item La FCEG approuve les nouveaux critères ajoutés par le BCAPG, car ces membres du corps professoral titulaires d'un permis d'exercice de l'ingénierie amélioreront la disponibilité et la qualité des possibilités d'apprentissage pratique offertes aux étudiants, ce qui préparera ces derniers à entrer sur le marché du travail après l'obtention de leur diplôme.
 \end{itemize}
\end{itemize}

\stanceheading{Les recommandations aux partenaires, aux parties intéressées et aux autres entités}
\begin{itemize}
 \item La FCEG demande au BCAPG de mettre à jour ses pratiques en matière de visites d'accréditation afin d'intégrer et de recueillir davantage de commentaires de la part des étudiants concernant l'apprentissage pratique, y compris les aspects pratiques des cours de base.
 \item La FCEG demande au BCAPG d'établir davantage de spécifications quant à la quantité d'apprentissage pratique requise pour l'obtention d'un diplôme d'ingénieur.
 \item La FCEG demande aux DDIC de recommander à leurs facultés membres de mieux construire les futurs programmes d'études afin d'obtenir un équilibre entre la théorie et les compétences pratiques dans les cours de base que les étudiants en ingénierie sont tenus de suivre.
 \item La FCEG demande aux DDIC de demander à ses facultés membres de fournir divers soutiens, tels que des clubs d'étudiants, afin de promouvoir les possibilités d'apprentissage pratique avec le corps enseignant.
 \item La FCEG demande en outre à ses facultés membres de compiler et de résumer le large éventail de possibilités d'apprentissage pratique offertes à leurs étudiants par le biais d'une page de ressources unifiée, telle qu'un carrefour de l'apprentissage expérientiel. Cette initiative permettra de s'assurer que tous les étudiants en ingénierie peuvent facilement découvrir et accéder aux possibilités d'apprentissage pratique pour enrichir leur diplôme.
 \item La FCEG demande aux DDIC de demander à ses facultés membres de recueillir les commentaires des étudiants concernant les possibilités d'apprentissage pratique.
 \item La FCEG demande à ses organisations régionales partenaires (WESST, CAÉGO, CRÉIQ, CÉGA) de collaborer à des efforts de sensibilisation afin de garantir un équilibre entre la théorie et les compétences pratiques dans leurs régions respectives.
 \item La FCEG invite les associations d'étudiants en génie à soutenir les efforts de promotion nationaux et à tirer parti de la mise à jour de la politique du BCAPG pour promouvoir la mise en œuvre de l'apprentissage pratique dans leur programme.
\end{itemize}

\stanceheading{Ce que fait la FCEG}
\begin{itemize}
 \item La FCEG offre actuellement la possibilité de participer aux cours du \textit{Board of European Students of Technology} (BEST) dans le cadre d'un partenariat, ce qui permet aux étudiants d'aborder un sujet technique qui pourrait être lié aux compétences professionnelles.
 \item La FCEG organise chaque année la Compétition canadienne d'ingénierie (CCI), qui intègre l'innovation, la créativité, la théorie et les compétences pratiques dans ses huit (8) volets de compétition : Réingénierie, débat oratoire, conception d'équipe junior, conception d'équipe senior, programmation, conception innovatrice, génie conseil et communications en génie.
 \item La FCEG organise des concours de cas lors de conférences sélectionnées, offrant aux étudiants une plateforme pour affiner et mettre en valeur leurs compétences en matière de résolution de problèmes et appliquer leurs connaissances académiques.
\end{itemize}

\stanceheading{Ce que la FCEG envisage de faire}
\begin{itemize}
 \item La FCEG recueillera des données sur l'apprentissage pratique dans le cadre de l'enquête nationale.
 \item La FCEG assurera le suivi avec les parties prenantes concernées, y compris les DDIC, afin d'aborder les modifications à apporter aux programmes d'études pour accroître les possibilités d'apprentissage pratique dans les facultés d'ingénierie canadiennes.
 \item La FCEG soulignera l'importance et la nécessité de l'apprentissage pratique lors des réunions du conseil d'administration du BCAPG et d'Ingénieurs Canada.
 \item La FCEG maintiendra des communications continues avec la \textit{International Federation of Engineering Education Societies} (\textit{IFEES}) concernant les moyens de mieux défendre les possibilités d'apprentissage pratique et les programmes d'échange offerts aux étudiants canadiens en génie.
 \item La FCEG collaborera avec l'Association canadienne de l'éducation en génie (ACEG) pour améliorer la qualité et la quantité des possibilités d'apprentissage pratique dans les écoles membres.
\end{itemize}

\newpage
